\documentclass{article}
\usepackage{lmodern}
\usepackage{amsmath}
\usepackage{listings}
\usepackage{fullpage}
\usepackage{parskip}
\usepackage{xcolor}
\lstset{
	basicstyle=\ttfamily,
	columns=fullflexible,
	frame=single,
	breaklines=true,
	postbreak=\mbox{\textcolor{red}{$\hookrightarrow$}\space},
}
\begin{document}
\title{Experiment `p\_5\_3\_first\_digit' Results}
\author{\today}
\date{}
\maketitle
\textbf{Experiment outcome: }SUCCESS
\\\textbf{Bad responses: }0
\\\textbf{Responses containing}~\texttt{assume}~\textbf{: }0
\\\textbf{Responses containing}~\texttt{expect}~\textbf{: }0
\\\textbf{Resolution attempts: }1
\\\textbf{Hard fails (resolution): }0
\\\textbf{Soft fails (resolution): }0
\\\textbf{Verification attempts: }1
\section*{Problem Specification}
\textbf{Problem name: }p\_5\_3\_first\_digit
\\\textbf{Natural language statement: }null
\\\textbf{Method signature: }p\_5\_3\_first\_digit(n: int) returns (first: int)
\subsection*{Ensures}
\begin{itemize}
\item \texttt{0 <= first < 10}
\item \texttt{intToSeq(n)[0] == first}
\end{itemize}
\subsection*{Requires}
\begin{itemize}
\item \texttt{0 <= n}
\end{itemize}
\subsection*{Functional Code Given}
\begin{lstlisting}
function intToSeq(n: int): seq<int>
  requires n >= 0
  ensures forall d :: d in intToSeq(n) ==> 0 <= d < 10
{
  if n < 10 then [n]
  else intToSeq(n / 10) + [n % 10]
}
\end{lstlisting}
\clearpage
\section*{GenAI interactions}
Below you will find all interactions between the `user' (program) and the `assistant' (OpenAI).
\subsection*{Program $\rightarrow$ GenAI}
\begin{lstlisting}
You are given the following task to perform in Dafny:

Write a method to satisfy the below requirements.

The signature should be:

p_5_3_first_digit(n: int) returns (first: int)

The method should respect the following contract:

requires 0 <= n, ensures 0 <= first < 10, ensures intToSeq(n)[0] == first

The contract uses the following dafny code:

function intToSeq(n: int): seq<int>
  requires n >= 0
  ensures forall d :: d in intToSeq(n) ==> 0 <= d < 10
{
  if n < 10 then [n]
  else intToSeq(n / 10) + [n % 10]
}

These function(s) must not be used in your implementation of the method.

Produce and show only the Dafny body of this method, including the curly braces that surround it. Do not show the signature nor contract. You must not use 'assume' nor 'expect' anywhere in your code.
\end{lstlisting}
\subsection*{GenAI $\rightarrow$ Program}
\textbf{System fingerprint: }null
\\\textbf{ID: }chatcmpl-DnQu9r22AVSoYJLhBxNDD75tf4fyH
\\\textbf{Model: }gpt-5.5-2026-04-23
\\\textbf{Created at: }1780673865
\\\textbf{Finish reason: }stop
\\\textbf{Usage: }464 tokens in, and 570 tokens out
\begin{lstlisting}
{
  if n < 10 {
    first := n;
  } else {
    first := p_5_3_first_digit(n / 10);
  }
}
\end{lstlisting}
\section*{Final Program}
\begin{lstlisting}
function intToSeq(n: int): seq<int>
  requires n >= 0
  ensures forall d :: d in intToSeq(n) ==> 0 <= d < 10
{
  if n < 10 then [n]
  else intToSeq(n / 10) + [n % 10]
}

method p_5_3_first_digit(n: int) returns (first: int)
	requires 0 <= n
	ensures 0 <= first < 10
	ensures intToSeq(n)[0] == first
{
  if n < 10 {
    first := n;
  } else {
    first := p_5_3_first_digit(n / 10);
  }
}
\end{lstlisting}
\section*{Total Token Usage}
\textbf{Input tokens: }464
\\\textbf{Output tokens: }570
\\\textbf{Reasoning tokens: }512
\\\textbf{Sum of `total tokens': }1034
\section*{Experiment Timings}
\textbf{Overall Experiment} started at 1780673864366, ended at 1780673904559, lasting 40193ms (40.19 seconds)
\\\textbf{Iteration \#1} started at 1780673864382, ended at 1780673904556, lasting 40174ms (40.17 seconds)
\end{document}
